%!TEX encoding = UTF8
%!TEX root = 0-notes.tex

\chapter{Racines}

\dfn{racine}{
	Soit $f\in\K[x]$ un polynôme non constant à coefficients dans $\K$.
	Le nombre $r\in\K$ est appelé \emphindex{racine} ou \emphindex{zéro} de $f$ dès que
		\[ f(r) = 0. \]
}{dfn:racine}

% taken from exe:root-factor
% maybe talk about roots before... need refactor.
\exe{}{
	Soit $f\in\K[x]$ un polynôme non constant.
	Montrer que si $r\in\K$ est une racine de $f$ si et seulement si  $x-r$ divise $f$ dans $\K[x]$.
}{exe:root-factor2}{
	todo
}

\exe{}{
	Est-ce que $x^3 - 2x + x -2$ est divisible par $x-2$ dans $\Z[x]$ ?
	
	Est-ce que $x^2 -4x + 3$ est divisible par $x+1$ dans $\Q[x]$ ?
}{exe:divisible-root}{
	todo
}

\exe{}{
	Soit $f\in\K[x]$.
	Montrer que si $a|f$ dans $\K[x]$, alors les racines de $a$ sont des racines de $f$.
}{exe:root-factor3}{
	todo
}

\section{Théorème fondamental de l'algèbre}

% thm fondamental
\thm{fondamental de l'algèbre, de d'Alembert}{
	Soit $f\in\C[x]$ un polynôme non constant à coefficients complexes.
	Alors $f$ admet une racine dans $\C$.
}{thm:fond-alg}

\cor{}{
	Soit $f\in\C[x]$ un polynôme à coefficients complexes de coefficient dominant $a\in\C$.
	Alors $f$ se décompose comme produit de monômes
		\[ f(x) = a\prod_{z | f(z) = 0} (x-z)^{m(z)}, \]
	où $m(z)\in\N$ est la \emphindex{multiplicité} de la racine $z$ dans $f$.
}{cor:decomp-C}

\dfn{algébriquement clos}{
	Le corps $\C$ est \emphindex{algébriquement clos} car chaque polynôme non constant de $\C[x]$ admet une racine dans $\C$.
	De façon équivalente, chaque polynôme de $\C[x]$ se scinde en produit de monômes.
}{dfn:alg-clos}

\thm{clôture algébrique}{
	Pour tout corps $\K$, il existe un sur-corps $\overline{\K}$ qui est algébriquement clos.
}{thm:cloture-algebrique}
	
\nt{
	Une démonstration constructive du théorème \ref{thm:cloture-algebrique} fonctionne comme suit :
	pour chaque polynôme de $\K[x]$, ajouter une racine au corps $\K$, et continuer le processus.
	La construction de la clôture de $\Q$ peut donc se faire en raisonnant par récurrence car l'ensemble $\Q[x]$ est dénombrable comme union dénombrable d'ensembles dénombrables.
	
	En revanche, le nombre de polynômes de $\K[x]$ étant indénombrable en général, une récurrence classique ne suffit pas, il faut une \emphindex{récurrence transfinie}.
	Une démonstration du théorème \ref{thm:cloture-algebrique} nécessite donc l'axiome du choix.
}


\exe{difficulty=1}{
	Soient $a, f\in\K[x]$ et $\overline{\K}$ la clôture algébrique de $\K$.
	Montrer que $a|f$ dans $\K[x]$ si et seulement si pour chaque racine de $a$, sa multiplicité dans $a$ est inférieure ou égale à celle dans $f$.
}{exe:root-factor4}{
	todo
}

\exe{difficulty=1}{
	Soit $f\in\K[x]$ et $f'$ sa dérivée formelle.
	Montrer que $\pgcd(f,f') = 1$ si et seulement si $f$ n'admet que des racines simples.
}{exe:gcd-simple-roots}{
	todo
}

\section{L'anneau $\R[x]$}

\lem{}{
	Soit $f\in\R[x]$ un polynôme à coefficients réels.
	Si $z\in\C\setminus\R$ est une racine de $f$, alors $\overline{z}$ aussi, et le polynôme 
		\[ (x - z)(x-\overline{z}) \]
	divise $f$ dans $\R[x]$.
}{lem:racines-conjuguees}

\pf{}{
 	D'abord, comme $f$ est à coefficients réels, $0 = \overline{f(z)} = f\left(\overline{z}\right)$, et $\overline{z}$ est bien une racine de $f$.
 	De plus, $x-z$ et $x-\overline{z}$ sont premiers entre eux sur $\C[x]$ car $z \neq \overline{z}$.
 	Il suit donc que $(x-z)(x-\overline{z})$ divise $f$ dans $\C[x]$.
 	
 	Or, $p(x) = (x-z)(x-\overline{z}) = x^2 - (z+\overline{z}) x + |z|^2$ est un polynôme à coefficients réels.
 	Montrons que $p$ divise $f$ dans $\R[x]$ également.
 	Considérons la division euclidienne $f = pq + r, \deg(r) < \deg(p)=2$ dans $\R[x]$.
 	Vu comme élément de $\C[x]$, le reste $r = f-pq$ s'annule en $z$ et en $\overline{z}$, il est donc divisible par $p$.
 	Comme son degré est strictement inférieur à celui de $p$, le reste $r$ doit être nul.
}

\cor{}{
	Soit $f\in\R[x]$ un polynômes à coefficients réels.
	D'après le théorème \ref{thm:decomp-Kx}, $f$ se décompose dans $\R[x]$ de façon unique comme produit de facteurs irréductibles.
	
	Chaque facteur est de degré 1 ou 2.
}{cor:decomp-Rx}

\exe{}{
	Montrer le corollaire \ref{cor:decomp-Rx}.
}{exe:cor:decomp-Rx}{
	todo
}

\section{Relations entre racines et coefficients}

\prop{}{
	Soit $\Lambda$ un ensemble fini d'éléments de $\K$.
	Considérons le produit 
		\[ f(x) = \prod_{\lambda \in \Lambda} (x+\lambda), \]
	polynôme de $\K[x]$.
	Alors $f(x)$ se développe en prenant le produit de soit $x$, soit $\lambda$ dans chaque monôme $x+\lambda$, et en sommant ces produits.
	
		\[  \prod_{\lambda \in \Lambda} (x+\lambda) = \sum_{k=0}^{|\Lambda|} x^k \sum_{\substack{I \subset \Lambda, \\ |I| = |\Lambda|-k}} \prod_{\lambda \in I} \lambda. \]
}{prop:expand-product}

\ex{}{
	Considérons le polynôme de $\R[x]$
		\[ f(x) = (x+2)(x+6)(x-1) \]
	qu'on souhaite développer.
	
	\begin{enumerate}[leftmargin=100pt, label=\bf \underline{Terme en $x^{\arabic*}$ :}, start=0]
		\item
		Pour obtenir un terme sans $x$, il faut choisir la constante dans chaque terme du produit.
		Celle-ci est donc $2\times6\times(-1) = -12$.
		Nous pouvons nous rassurer que c'est vrai en calculant $f(0)$.
		
		\item
		Pour obtenir un terme en $x$, il faut choisir un seul $x$ parmis les trois facteurs.
		Les trois possibilités sont $x\times6\times(-1)$, $2\times x\times (-1)$, et $2\times6\times x$, qu'on somme pour obtenir $-6x-2x+12x = 4x$.
		
		\item
		Pour obtenir un terme en $x^2$, il faut choisir deux $x$ parmis les trois facteurs.
		Les trois possibilités sont $2\times x \times x$, $x \times 6 \times x$, et $x \times x \times (-1)$, qu'on somme pour obtenir
		$2x^2 + 6x^2 - x^2 = 7x^2$.
		
		\item
		Pour obtenir un terme en $x^3$, il faut choisir trois $x$, et la seule façon de faire est bien sûr $x\times x \times x = x^3$.
	\end{enumerate}
	
	Par conséquent, 
		\[ (x+2)(x+6)(x-1) = x^3 + 7x^2 + 4x - 12. \]
}{exe:dev-deg-3}


\exe{}{
	Soit $n, k\in\N$ deux entiers avec $k\leq n$. 
	Rappelons la notation ${n \choose k}$, lue « $k$ parmi $n$ », donnant le nombre de façons de choisir $k$ objets parmi $n$ objets indistinguables.
	
	Montrer que
		\[ (1+x)^n = \sum_{k=0}^n {n \choose k} x^k. \]
}{exe:binom-Newton}{
	todo
}

\exe{difficulty=1}{
	Soit $n\in\N$. En comparant les termes en $x^n$ de l'identité
		$(1+x)^{2n} = (1+x)^n \cdot (1+x)^n, $
	montrer que
		\[ \sum_{k=0}^{n} {n \choose k} {n \choose n-k} = {2n \choose n}. \]
}{exe:id-newton}{
	todo
}

\cor{}{
	Soit $f\in\K[x]$ un polynôme unitaire de degré $n$ se scindant sur $\K$ :
		\[ f(x) = \prod_{\lambda \in \Lambda} (x-\lambda), \]
	avec $\Lambda$ un $n$-tuple des racines de $f$ comptées avec multiplicités.
	
	Alors la forme développée de $f$ vérifie
		\[ f(x) = a_n x^n + a_{n-1}x^{n-1} + \hdots + a_1 x + a_0, \]
	où
	\begin{enumerate}
		\item $a_n = 1$ car $f$ est unitaire ;
		\item $a_0 = (-1)^n  \prod_{\lambda \in \Lambda} \lambda$ le produit des racines ;
		\item $a_{n-1} = -  \sum_{\lambda \in \Lambda} \lambda$ la somme des racines.
	\end{enumerate}
}{cor:rel-racine-coeff}

\exe{}{
	Dans les notations du corollaire \ref{cor:rel-racine-coeff},
	exprimer $a_2$ et $a_{n-2}$ en fonction des racines $\lambda$.
}{exe:a2}{
	todo
}

\exe{}{
	L'objectif de l'exercice est de démontrer le \emphindex{théorème de Wilson} qui énonce que $(n-1)! = -1 \mod n$ si et seulement si $n$ est premier.
	\begin{enumerate}
		\item
		Montrer que si $n$ n'est pas premier, alors $(n-1)! = 0 \mod n$.
		\item
		Supposons $n$ premier.
		Montrer que le polynôme $x^n - x$ admet $n$ racines distinctes dans $\F_n[x]$.
		\item
		Conclure en étudiant le coefficient constant de $x^{n-1} - 1$.
	\end{enumerate}
}{exe:Wilson}{
	todo
}

\exe{difficulty=2}{
	Soit $a, b\in\N^*$ deux entiers naturels tels que $ab+1$ divise $a^2 + b^2$.
	Montrer que $\frac{a^2 +b^2}{ab+1}$ est un carré parfait.
}{exe:Vieta-jump}{
	IMO 1988.
}


\exe{difficulty=1}{
	Soient $\Lambda \in \K^n$ \underline{toutes} les valeurs propres d'une matrice $A \in \K^{n\times n}$, comptées avec multiplicités.
	Ainsi, le polynôme caractéristique
	$\det(x I - A) = x^n + \alpha_{n-1} x^{n-1} + \dots + \alpha_1 x + \alpha_0$
	est unitaire et se scinde comme produit
		\[\det(x I - A) = \prod_{\lambda\in\Lambda}(x - \lambda). \]
	\begin{enumerate}
		\item
		Montrer d'abord que le déterminant d'une matrice est égal au produit de ses valeurs propres
		 	\[ \det(A) = \prod_{\lambda \in \Lambda} \lambda. \]
	\end{enumerate}
	À l'aide de la formule de Leibniz, 
		\[ \det(M) = \sum_{\sigma \in S_n} \sgn(\sigma) \prod_{i=1}^{n} M_{i, \sigma(i)}, \]
	montrer ensuite que
	\begin{enumerate}[resume]	
		\item $\alpha_{n-1} = (-1)^{n-1} \tr (A)$ ; et que
		\item La trace d'une matrice est égale à la somme de ses valeurs propres
			\[ \tr(A) = \sum_{\lambda \in \Lambda} \lambda. \]
	\end{enumerate}
}{exe:det-trace-val-p}{
	todo
}

\section{Racines $n$-ièmes de l'unité}

\nt{
	Soit $n\in\N^*$.
	Le polynôme $x^n - 1 \in\C[x]$ admet pour dérivée formelle $nx^{n-1}$, dont la seule racine est 0.
	Ainsi, $x^n -1$ est premier avec son polynôme dérivé, et il admet donc $n$ racines complexes distinctes.
}

\dfn{racine $n$-ième de l'unité}{
	Soit $z\in\C$ un nombre complexe et $n\in\N^*$.
	Le nombre $z$ est \emphindex{racine $n$-ième de l'unité} dès que $z^n = 1$.
}{dfn:racine-unite}

\notation{
	L'ensemble des nombres complexes de module 1 est noté	
		\[ \U = \bigset{ z\in\C \tq |z| = 1}. \]
	
	Pour $n\in\N^*$, l'ensemble des nombres racines $n$-ièmes de l'unité est noté 
		\[ \U_n = \bigset{ z \in \C \tq z^n = 1 }. \]
}

\exe{}{
	Soit $z\in\C$.
	Montrer l'implication suivante. La réciproque est-elle vraie ?
	\begin{center}
	\begin{tabular}{ccc}
		$\Bigpar{ \exists n\in\N^* \tq z\in\U_n }$ &  $\implies$ & $\Bigpar{ |z| = 1}$
	\end{tabular}
	\end{center}
}{exe:root-is-unit}{
	todo
}


\exe{difficulty=1}{
	Soit $n\in\N$ impair et $\zeta = e^{\frac{2 \pi i}n}$.
	Montrer l'identité suivante dans $\C[x,y]$.
		\[ x^n + y^n = \prod_{k=0}^{n-1} \bigpar{x + \zeta^k y} \]
	Est-elle vraie pour $n\in\N^*$ pair ?
}{exe:split-xn-p-yn}{
	todo
}

\exe{}{
	Soit $n\in\N^*$.
	Montrer que $\bigpar{\U_n, \times}$ est un groupe isomorphe à $\bigpar{ \Z/n\Z, +}$ en explicitant l'isomorphisme.
}{exe:Un-sous-groupe}{
	todo
}


\exe{}{
	Résoudre le système suivant pour $a, b, c \in\C$.
	\[\begin{cases*}
		abc = 1,  \\
		a+b+c = 0, \\
		ab+ac+bc = 0. \\
	\end{cases*}\]
}{exe:systeme-units}{

}

\section*{Exercices supplémentaires}
\addcontentsline{toc}{section}{Exercices supplémentaires}

\exe{}{
	Soit $z\in\C$ une racine $p$-ième de l'unité où $p$ est un nombre premier.
	Montrer que le polynôme minimal annulateur de $z$ est $\frac{x^p - 1}{x-1}$.
}{exe:pmin-unit-root}{
	todo
}

\exe{}{
	Soit $n\in\N^*$.
	En étudiant $(1+x)^n + (1-x)^n$, montrer que
		\[ \sum_{\substack{k=0 \\ \text{$k$ pair}}}^n {n \choose k} =  \sum_{\substack{k=0 \\ \text{$k$ impair}}}^n {n \choose k} =  2^{n-1}. \]
	Quelle partie du raisonnement échoue lorsque $n=0$ ?
}{exe:Newton-even-odd}{
	Les termes pairs sont doublés et les termes impairs disparaissent.
	Évaluer en $x=1$ conclut pour la première somme.
	La deuxième somme est déduite de la première en utilisant que $2^n - 2^{n-1} = 2^{n-1}$.

	Lorsque $n=0$, l'identité $(1-x)^0 = 1$ semble vraie mais elle ne l'est pas précisément en $x=1$.
}

\exe{}{
	Montrer que $\U$ est un sous-groupe de $\bigpar{\C^*, \times}$.
}{exe:U-sous-groupe}{
	todo
}

\exe{}{
	Soit $n\in\N^*$.
	Montrer que $\bigpar{\U_n, \times}$ est le seul sous-groupe d'ordre $n$ de $\U$.
}{exe:Un-sous-groupe-unique}{
	todo
}

\exe{}{
	Soit $n\in\N^*$.
	Décrire explicitement l'ensemble $\mu_n$ des éléments d'ordre $n$ du groupe $\U$.
}{exe:U-ordre-n}{
	todo
}

\exe{difficulty=1}{
	Soit $p>3$ un nombre premier. Montrer l'égalité suivante.
		\[ \prod_{1 \leq i < j \leq p-1} ij = 0 \mod p. \]
	Vérifier l'égalité pour $p=5$ et montrer qu'elle n'est pas vraie pour $p=3$.
}{exe:Wilson-bis}{
	En reprenant le raisonnement de l'exercice \ref{exe:Wilson} et le résultat de l'exercice \ref{exe:a2}, étudions le coefficient en $x^{p-2}$ du polynôme $x^{p-1} -1$ qui se scinde comme suit sur $\F_p[x]$.
		\[ x^{p-1} - 1 = \prod_{i=1}^{p-1} (x- i) \]
	Comme $p>3$, le coefficient en $x^{p-2}$ de $x^{p-1}-1$ est nul, et donc
		\[ \prod_{1 \leq i < j \leq p-1} (-i)(-j) = 0 \]
	dans $\F_p$.
}


\exe{difficulty=2}{
	Soit $n\in\N^*$.
	Pour $d\in\N, d|n$ diviseur de $n$, posons
		\begin{align*}
		\mu_d = \bigset{ e^{2\pi i \frac{k}d} \tq 1 \leq k \leq d \et \pgcd(k, d) = 1} && \et && \phi_d = \prod_{z \in \mu_d} \bigpar{ x-z }.
		\end{align*}
	\begin{enumerate}
		\item
		Montrer que 
			\[ x^n - 1 = \prod_{d|n} \phi_d(x). \]
		\item
		En déduire que si $p$ est premier, alors $\phi_p(x) = \frac{x^p - 1}{x-1}$ et qu'il est irréductible dans $\Z[x]$.
		\item
		Montrer que si $n$ est \emphindex{semi-premier}, c'est-à-dire que $n=pq$ avec $p, q$ premiers, alors
			\[ \phi_n(x) = \frac{x^n-1}{\phi_p(x) \phi_q(x) (x-1)} \]
		est une division dans $\C[x]$ et donc dans $\Q[x]$. Est-ce une division dans $\Z[x]$ ? Justifier.
		\item
		Généraliser pour montrer que $\phi_n \in \Z[x]$ pour tout $n\in\N^*$.
	\end{enumerate}
}{exe:phin-integer}{
	todo
}


\shipoutExercise

\section*{Solutions}
\addcontentsline{toc}{section}{Solutions}

\shipoutAnswer



